\documentclass[11pt]{article}

\usepackage[margin=1in]{geometry}
\usepackage{amsmath,amssymb}
\usepackage{array}
\usepackage{booktabs}
\usepackage{enumitem}
\usepackage{xcolor}
\usepackage{hyperref}
\usepackage{setspace}

\setstretch{1.10}
\setlength{\emergencystretch}{2em}

\hypersetup{
    colorlinks=true,
    linkcolor=blue!50!black,
    urlcolor=blue!50!black,
    citecolor=black
}

% --------------------
% Course information
% --------------------
\newcommand{\course}{DSA4213: Natural Language Processing for Data Science}
\newcommand{\semester}{Semester 1, AY2026/2027}

\begin{document}

\begin{center}
{\LARGE\bfseries Group Project Final Report}\\[0.4em]
{\large Project Title}\\[0.9em]
Group ID: \underline{\hspace{3cm}}\\[0.9em]
\course\\
\semester
\end{center}

% ====================
% Report structure, format, and length (delete this box before submission)
% ====================
\noindent\fbox{\parbox{\dimexpr\linewidth-2\fboxsep-2\fboxrule\relax}{%
\textbf{Report structure, format, and length} (delete this box before submission)
\begin{itemize}[leftmargin=*,itemsep=0.2em,topsep=0.3em]
    \item The section structure is a guide. You may combine, rename, or reorder
    sections to suit your project, as long as you address the relevant assessment
    criteria and include the required contribution statement.
    \item The main text must not exceed \textbf{8 pages}. The main text is everything
    before the Contribution Statement, including the title block, figures, and tables.
    The limit is strict. Reports whose main text exceeds 8 pages will be penalized.
    \item The Contribution Statement, References, and Appendix do not count toward the
    page limit.
    \item Keep the 11pt font, 1-inch margins, and line spacing set in this template.
    Reports that change these settings will be treated as exceeding the page limit.
    \item The main text must be self-contained. The Appendix may contain supplementary
    results, prompts, or agent traces, but graders are not required to read it.
\end{itemize}}}

% ====================
% Administrative information
% ====================
\section*{Administrative Information}

\begin{tabular}{@{}>{\bfseries\raggedright\arraybackslash}p{0.25\linewidth}
                    >{\raggedright\arraybackslash}p{0.70\linewidth}@{}}
Group members & Name (Student ID) for each member. \\[0.7em]
Project direction & One short phrase describing the project area. \\[0.7em]
Repository & Link to the code and project materials: \url{https://} \\
\end{tabular}

% ====================
% Executive summary
% ====================
\section*{Project Summary (Optional)}

Summarize the agent-related question, what the group built or studied, how it was
evaluated, and the main findings, in one paragraph.

% ====================
% 1. Problem and motivation
% ====================
\section*{1. Problem, Research Question, and Motivation}

State the question, objective, or hypothesis that the project addresses. Explain the
practical or scientific motivation, the intended users or setting, and why an
\emph{agent-based} approach is relevant. Define what success means for this project.

If the question, scope, or design changed after the proposal, state the main change
and the reason in two or three sentences.

% ====================
% 2. Related work and contribution
% ====================
\section*{2. Related Work and Contribution}

Identify the prior systems, datasets, benchmarks, papers, or tools closest to the
project. State what the group contributed beyond assembling existing components.
Cite the models, APIs, datasets, environments, frameworks, and code that the project
uses.

% ====================
% 3. Agent setup
% ====================
\section*{3. Agent and Task Setup}

Describe the final setup precisely. For an empirical study, describe the agent
system and the study design or conditions considered. The list below is a guide,
not a checklist: you do not need to address every item. Include only the elements
that are relevant to your project:

\begin{itemize}[leftmargin=*,itemsep=0.25em]
    \item task and operating environment;
    \item inputs or observations and outputs or actions;
    \item model(s), prompts, tools, APIs, retrieval sources, memory, or other agents;
    \item users, simulators, datasets, benchmarks, or interaction protocols; and
    \item stopping conditions, constraints, permissions, and failure recovery.
\end{itemize}

% ====================
% 4. System or study design
% ====================
\section*{4. System or Study Design}

Describe what the group implemented or investigated and how the main components fit
together. Identify the part that is technically meaningful and owned by the group,
and explain the reasons for the main design choices. If useful, include a system
diagram, pseudocode, or an example agent trajectory.

Include the implementation details needed to understand the results. Further detail
can go in the repository or the Appendix.

% ====================
% 5. Evaluation
% ====================
\section*{5. Evaluation and Results}

Describe the evaluation protocol: the tasks or test cases, the metrics, rubrics, or
other evidence used, the number of runs or examples, and, where relevant, how
variability was handled. Then report the results and connect each main claim to the
evidence that supports it.

Present the results in whichever format best suits the project. Tables, figures,
and short case studies are all acceptable.

\subsection*{5.1 Comparisons and Ablations (if applicable)}

Report any baselines, controlled comparisons, or ablations that help answer the
project question, and explain what each one shows.

% ====================
% 6. Analysis
% ====================
\section*{6. Analysis of Results, Failures, Limitations, Costs, and Risks}

Interpret the results. Use representative successes and failures, such as examples
or agent traces, to explain the aggregate results. Describe the main failure modes
and, where possible, their causes.

Report the costs that matter for the project, such as API cost, tokens, latency, or
compute. State the main limitations and risks of the system or study. Well-supported
negative or mixed findings are acceptable.

% ====================
% 7. Reproducibility and responsible practice
% ====================
\section*{7. Reproducibility and Responsible Practice}

Explain how the main results can be reproduced from the submitted materials: where
the code and setup instructions are, which models and versions were used, and which
data or environments are required. Refer to the setup instructions in the repository
rather than repeating them here.

Describe any safeguards for privacy, security, permissions, or safe use. Do not
include credentials, private data, or material that the group does not have
permission to share.

% ====================
% 8. Conclusion
% ====================
\section*{8. Conclusion (Optional as a separate section)}

Summarize the main findings in a short paragraph.

% ====================
% End of main text (8-page limit)
% ====================
\clearpage

% ====================
% Contribution statement (not counted toward the page limit)
% ====================
\section*{Contribution Statement}

List each member's main contributions and an approximate contribution share. The
shares must sum to 100\%. If contributions were broadly comparable, equal shares are
appropriate. The group should review this statement together before submission.

\begin{center}
\begin{tabular}{@{}>{\raggedright\arraybackslash}p{0.28\linewidth}
                >{\raggedright\arraybackslash}p{0.52\linewidth}
                >{\centering\arraybackslash}p{0.12\linewidth}@{}}
\toprule
\textbf{Member (Student ID)} & \textbf{Main contributions} & \textbf{Share (\%)} \\
\midrule
Member 1 & & \\
Member 2 & & \\
Member 3 & & \\
Member 4 & & \\
Member 5 & & \\
\midrule
\textbf{Total} & & \textbf{100} \\
\bottomrule
\end{tabular}
\end{center}

% ====================
% References (not counted toward the page limit)
% ====================
\bibliographystyle{plain}
% \bibliography{references}

% ====================
% Appendix (optional, not counted toward the page limit)
% ====================
\appendix
\section*{Appendix (Optional)}

Supplementary results, prompts, additional examples, or agent traces. The main text
must be understandable without the Appendix.

\end{document}
